\documentclass[lualatex,ja=standard,a4paper,12pt]{bxjsarticle}
\usepackage{amsmath, amssymb, amsthm}
\usepackage{geometry}
\usepackage{xcolor}
\usepackage{tcolorbox}
\usepackage{booktabs}
\usepackage{array}
\usepackage{graphicx}
\usepackage{tikz}
\usetikzlibrary{arrows.meta, positioning}

\geometry{margin=25mm}

\tcbuselibrary{theorems, skins, breakable}

\newtcolorbox{defbox}[1]{
  colback=blue!5, colframe=blue!60!black,
  fonttitle=\bfseries, title={#1}, arc=3pt, breakable
}
\newtcolorbox{thmbox}[1]{
  colback=green!5, colframe=green!50!black,
  fonttitle=\bfseries, title={#1}, arc=3pt, breakable
}
\newtcolorbox{pointbox}{
  colback=orange!8, colframe=orange!70!black,
  fonttitle=\bfseries, title={ポイント}, arc=3pt, breakable
}
\newtcolorbox{exbox}[1]{
  colback=gray!6, colframe=gray!50,
  fonttitle=\bfseries, title={例：#1}, arc=3pt, breakable
}
\newtcolorbox{warningbox}{
  colback=red!5, colframe=red!60!black,
  fonttitle=\bfseries, title={注意}, arc=3pt, breakable
}

\title{{\large Hoff 第6章} \\[6pt]
       \textbf{6.5\; General Properties of the Gibbs Sampler} \\[4pt]
       {\normalsize （ギブスサンプラーの一般的性質）ゼミ資料}}
\author{}
\date{}

\begin{document}
\maketitle
\thispagestyle{empty}

% ===========================
\section{この節の位置づけ}
% ===========================

6.4節では，2パラメータの具体例でギブスサンプラーを走らせ，
結果が離散近似とほぼ一致することを見た．
しかしそれは\textbf{1つの例でうまくいった}というだけで，

\begin{itemize}
  \item なぜ適当な出発点から始めてよいのか
  \item パラメータが2個でなく $p$ 個でも同じことが言えるのか
  \item そもそも何が保証されているのか
\end{itemize}

\noindent には答えていない．本節はその整理にあてられる．

\medskip
内容は大きく2つに分かれる：
\textbf{(1) 一般形と収束の主張}，
\textbf{(2) 「推定」と「近似」を区別せよという概念的な注意}．
後半は数式こそ出てこないが，第7章以降をずっと支える考え方である．

% ===========================
\section{一般の $p$ パラメータへの拡張}
% ===========================

パラメータベクトルを $\phi = \{\phi_1, \ldots, \phi_p\}$ とし，
関心のある確率分布を $p(\phi) = p(\phi_1, \ldots, \phi_p)$ とする．
正規モデルなら $\phi = \{\theta, \sigma^2\}$，
$p(\phi)$ は事後分布 $p(\theta, \sigma^2 \mid y_1, \ldots, y_n)$ である．

\begin{defbox}{ギブスサンプラー（一般形）}
出発点 $\phi^{(0)} = \{\phi_1^{(0)}, \ldots, \phi_p^{(0)}\}$ を与え，
$\phi^{(s-1)}$ から $\phi^{(s)}$ を次のように生成する：
\begin{align*}
  &1.\quad \phi_1^{(s)} \sim p\!\left(\phi_1 \mid \phi_2^{(s-1)}, \phi_3^{(s-1)}, \ldots, \phi_p^{(s-1)}\right) \\
  &2.\quad \phi_2^{(s)} \sim p\!\left(\phi_2 \mid \phi_1^{(s)}, \phi_3^{(s-1)}, \ldots, \phi_p^{(s-1)}\right) \\
  &\qquad \vdots \\
  &p.\quad \phi_p^{(s)} \sim p\!\left(\phi_p \mid \phi_1^{(s)}, \phi_2^{(s)}, \ldots, \phi_{p-1}^{(s)}\right)
\end{align*}
\end{defbox}

\begin{pointbox}
\textbf{条件の中身が途中で入れ替わっていく}

$j$ 番目を引くときに条件とするのは，
\textbf{すでに更新済みの $\phi_1, \ldots, \phi_{j-1}$ は新しい値 $(s)$}，
\textbf{まだ更新していない $\phi_{j+1}, \ldots, \phi_p$ は古い値 $(s-1)$} である．

\medskip
1ステップ進むごとに「新しい値」の領域が左から右へ1つずつ伸びていき，
$p$ 番目まで終わると全部が新しい値になって $\phi^{(s)}$ が完成する．
2パラメータの場合（6.4節）に
$\sigma^2$ の更新で $\theta^{(s)}$（新しいほう）を使ったのと同じ規則である．
\end{pointbox}

% ===========================
\section{マルコフ連鎖であること}
% ===========================

この手続きが生成するのは，ベクトルの従属列である：
\[
  \phi^{(1)} = \{\phi_1^{(1)}, \ldots, \phi_p^{(1)}\}, \quad
  \phi^{(2)} = \{\phi_1^{(2)}, \ldots, \phi_p^{(2)}\}, \quad
  \ldots, \quad
  \phi^{(S)} = \{\phi_1^{(S)}, \ldots, \phi_p^{(S)}\}.
\]

\begin{defbox}{マルコフ性とマルコフ連鎖}
この列において，$\phi^{(s)}$ が
$\phi^{(0)}, \ldots, \phi^{(s-1)}$ に依存するのは
\textbf{$\phi^{(s-1)}$ を通じてのみ}である．
言い換えれば，$\phi^{(s-1)}$ を与えたとき
$\phi^{(s)}$ は $\phi^{(0)}, \ldots, \phi^{(s-2)}$ と条件付き独立である．

\medskip
この性質を\textbf{マルコフ性}（Markov property）といい，
これを持つ列を\textbf{マルコフ連鎖}（Markov chain）と呼ぶ．
\end{defbox}

\begin{pointbox}
\textbf{なぜマルコフ性が成り立つのか}

アルゴリズムを見れば明らかである：
$\phi^{(s)}$ を作るのに使う情報は，
$\phi^{(s-1)}$ の成分とデータだけであって，
それより前の履歴はどこにも現れない．

\medskip
これは「記憶がない」ということであり，
実装上は\textbf{直前の状態だけ保持すればよい}ことを意味する
（6.4節の R コードが \texttt{phi} という長さ2のベクトル1本を
使い回していたのはこのためである）．
\end{pointbox}

% ===========================
\section{収束についての2つの主張}
% ===========================

Hoff は「本書で扱うすべてのモデルで満たされるいくつかの条件のもとで」
と前置きしたうえで，次の2つを述べる．

\begin{thmbox}{主張1：分布の収束}
パラメータ空間の任意の集合 $A$ について
\[
  \Pr(\phi^{(s)} \in A) \;\longrightarrow\; \int_A p(\phi)\, d\phi
  \qquad (s \to \infty).
\]
すなわち，$\phi^{(s)}$ の\textbf{標本分布は目標分布に近づく}．
\textbf{出発点 $\phi^{(0)}$ が何であってもよい}
（ただし出発点によって，目標に達するまでの速さは違う）．
\end{thmbox}

\begin{thmbox}{主張2：平均の収束（式 6.1）}
関心のあるたいていの関数 $g$ について
\[
  \frac{1}{S}\sum_{s=1}^{S} g(\phi^{(s)})
  \;\longrightarrow\;
  E[g(\phi)] = \int g(\phi)\, p(\phi)\, d\phi
  \qquad (S \to \infty).
\]
\end{thmbox}

\begin{pointbox}
\textbf{6.3節の宿題が片づいた}

6.3節で示せたのは「同時事後分布から始めれば，そこから外れない」
（不動点であること）までだった．
\textbf{主張1はそれより強く，どこから始めても目標に近づく}と言っている．
これがあってはじめて，6.4節の「出発点は適当でよい」が正当化される．
\end{pointbox}

\begin{warningbox}
\textbf{「いくつかの条件」の中身は書かれていない}

Hoff は \textit{``Under some conditions that will be met for all of the
models discussed in this text''} と述べるだけで，条件を明示していない．

\medskip
教科書の範囲を出るが，大まかには
\textbf{連鎖がパラメータ空間のどこへでも到達できること}（既約性）と，
\textbf{周期的な振動に陥らないこと}（非周期性）が要る．
たとえば完全条件付き分布の台が分断されていて
連鎖が片側から出られないような場合には，主張1は成り立たない．

\medskip
なお，条件が満たされていても\textbf{収束が実用的な時間内に起きるとは限らない}．
$s \to \infty$ の「$\infty$」がどれほど遠いかは別問題であり，
それが6.6節の主題になる．
\end{warningbox}

\paragraph{MCMC という名前}
主張2は「標本平均で期待値を近似してよい」と言っており，
これは第4章の\textbf{モンテカルロ近似とまったく同じ使い方}である．
違いは，標本が独立ではなく\textbf{マルコフ連鎖}として生成される点だけ．
そこでこの種の近似を
\textbf{マルコフ連鎖モンテカルロ（Markov chain Monte Carlo, MCMC）近似}，
手続きを\textbf{MCMCアルゴリズム}と呼ぶ．

\begin{exbox}{ユスリカデータへの適用}
6.4節の1000個のギブス標本について，式(6.1)は
$\{(\theta^{(1)}, \sigma^{2(1)}), \ldots, (\theta^{(1000)}, \sigma^{2(1000)})\}$
の同時分布が $p(\theta, \sigma^2 \mid y_1, \ldots, y_n)$ に近いことを意味し，
\begin{align*}
  E[\theta \mid y_1, \ldots, y_n]
    &\;\approx\; \frac{1}{1000}\sum_{s=1}^{1000} \theta^{(s)} = 1.804, \\[4pt]
  \Pr\!\left(\theta \in [1.71,\ 1.90] \mid y_1, \ldots, y_n\right) &\approx 0.95.
\end{align*}
\end{exbox}

% ===========================
\section{「推定」と「近似」を区別する}
% ===========================

ここからは数式のない，しかし重要な話である．

\medskip
ベイズのデータ解析をモンテカルロ法で行っていると，
サンプリング手続きと確率分布が入り乱れて混乱しやすい．
そこで，解析のうち\textbf{統計的な部分}と
\textbf{数値的な近似の部分}を分けて考えることが助けになる．

\begin{thmbox}{ベイズのデータ解析の3要素（第1章）}
\begin{enumerate}
  \item \textbf{モデルの指定}：データの標本分布を表すと考えられる
        確率分布の族 $\{p(y \mid \phi),\ \phi \in \Phi\}$

  \item \textbf{事前分布の指定}：どのパラメータ値がもっともらしいかについての
        事前情報を表す確率分布 $p(\phi)$

  \item \textbf{事後分布の要約}：事後分布 $p(\phi \mid y)$ を，
        事後平均・中央値・最頻値・予測確率・信用域といった
        関心のある量で記述すること
\end{enumerate}
\end{thmbox}

\noindent 1と2を決め，データを集めた時点で，事後分布は
\[
  p(\phi \mid y) = \frac{p(\phi)\,p(y \mid \phi)}{p(y)}
  = \frac{p(\phi)\,p(y \mid \phi)}{\int p(\phi)\,p(y \mid \phi)\, d\phi}
\]
として\textbf{完全に決まってしまう}．
ある意味で，そこから先にモデリングも推定も残っていない．
残っているのは3の「要約」だけである．

\medskip
ところが多くのモデルで $p(\phi \mid y)$ は複雑で書き下すのも難しい．
そこで，$p(\phi \mid y)$ を「\textbf{眺める}」ための便利な方法として
そこからのモンテカルロ標本を調べる．

\begin{warningbox}
\textbf{モンテカルロ法・MCMC法は何でないか}

\begin{itemize}
  \item \textbf{モデルではない}
  \item $y$ と $p(\phi)$ に含まれる以上の
        \textbf{「情報」を生み出しはしない}
  \item 単に $p(\phi \mid y)$ を\textbf{「眺める方法」にすぎない}
\end{itemize}

\medskip
たとえば $\phi^{(1)}, \ldots, \phi^{(S)}$ が $p(\phi \mid y)$ からの
近似的な標本であるとき，これらは事後分布を記述するのに役立つ：
\begin{align*}
  \frac{1}{S}\sum_s \phi^{(s)} &\approx \int \phi\, p(\phi \mid y)\, d\phi, \\[4pt]
  \frac{1}{S}\sum_s \mathbf{1}(\phi^{(s)} \le c)
    &\approx \Pr(\phi \le c \mid y) = \int_{-\infty}^{c} p(\phi \mid y)\, d\phi.
\end{align*}
\end{warningbox}

\begin{defbox}{用語の使い分け}
この区別を忘れないために，次のように語を使い分けるとよい：
\begin{itemize}
  \item \textbf{推定}（estimation）：
        $p(\phi \mid y)$ を使って $\phi$ について推論すること
  \item \textbf{近似}（approximation）：
        モンテカルロ的な手続きで積分を近似すること
\end{itemize}
\end{defbox}

\begin{pointbox}
\textbf{なぜこの区別が大事か}

標本サイズ $S$ を増やして改善されるのは\textbf{近似}のほうだけである．
$S$ を 1000 から 100{,}000 にすれば事後平均の数値は精密になるが，
\textbf{事後分布そのものは1ミリも変わらない}．
$p(\phi \mid y)$ はデータと事前分布だけで決まっているからである．

\medskip
逆に，事後分布に不満があるなら（たとえば区間が広すぎるなら），
直すべきなのはモデルか事前分布かデータであって，
MCMCの反復回数ではない．

\medskip
また6.6節で出てくる\textbf{MCMC標準誤差}は，
「$\phi$ がどれだけ不確実か」ではなく
「\textbf{われわれの数値計算がどれだけ雑か}」を測る量である．
事後標準偏差とMCMC標準誤差を混同しないこと．
\end{pointbox}

% ===========================
\section{まとめ}
% ===========================

\begin{thmbox}{6.5節のまとめ}
\begin{enumerate}
  \item ギブスサンプラーは $p$ 個のパラメータに自然に拡張される．
        $j$ 番目を引くときは，更新済みの成分には新しい値を，
        未更新の成分には古い値を使う．

  \item 生成される列は\textbf{マルコフ連鎖}である
        （$\phi^{(s)}$ は $\phi^{(s-1)}$ を通じてのみ過去に依存する）．

  \item ある条件のもとで，
        \textbf{出発点によらず}標本分布は目標分布に収束し（主張1），
        標本平均は期待値に収束する（主張2，式6.1）．
        これで6.3節の循環が正式に解消される．
        この種の近似を\textbf{MCMC近似}と呼ぶ．

  \item モンテカルロ法とMCMC法は\textbf{モデルでも情報源でもなく}，
        すでに決まっている事後分布を「眺める方法」にすぎない．
        $p(\phi\mid y)$ による推論を\textbf{推定}，
        積分の数値計算を\textbf{近似}と呼び分けて区別する．
\end{enumerate}
\end{thmbox}

\paragraph{次節へ}
本節の主張1・2はいずれも $s \to \infty$，$S \to \infty$ の話である．
現実には有限回しか回せない．

\medskip
\textbf{有限の $S$ で，その近似はどれくらい信用できるのか？}
\textbf{収束したかどうかをどう判断するのか？}
——これが6.6節の問題である．

\end{document}
